% ------------------------- Prefacio ---------------------------------
\chapter*{Prefácio}
\addcontentsline{toc}{chapter}{Prefácio}
\markboth{Prefácio}{Prefácio}

A modelagem computacional de \qubit{}s baseados em pontos quânticos deixou de
ser uma curiosidade acadêmica para se tornar uma competência estratégica.
Governos e empresas investem bilhões em computação quântica, e há uma demanda
crescente por profissionais capazes de descrever, com rigor atômico, os
dispositivos que hão de sustentar essa tecnologia. No entanto, quem tenta
ingressar na área costuma esbarrar em dois obstáculos: pacotes de estrutura
eletrônica com curvas de aprendizado íngremes e, muitas vezes, com custos de
licenciamento proibitivos.

O \pyscf{} --- \emph{Python-based Simulations of Chemistry Framework} --- derruba
essas barreiras. É livre, gratuito e integralmente programável em \python{}, a
linguagem que se tornou a \emph{lingua franca} da ciência computacional. Ainda
assim, mesmo uma ferramenta amigável pode intimidar quem nunca executou um
cálculo de química quântica ou escreveu uma linha de código.

Este livro nasce de uma convicção simples: \textbf{é possível partir do zero
absoluto e chegar à simulação completa de um \qubit{} real}. Nós o faremos com
uma abordagem prática, na qual cada conceito é apresentado por meio de exemplos
executáveis. Mais do que isso, todo o material converge para um \textbf{projeto
longitudinal único} --- a construção, passo a passo, de um modelo de \qubit{} de
spin em silício dopado com fósforo (Si:P). Você começará digitando
\code{import pyscf} e terminará extraindo o acoplamento hiperfino, o tensor
\emph{g}, a energia de troca entre \qubit{}s e estimativas de tempos de
descoerência.

Este primeiro volume cobre a \textbf{Parte 0 --- Preparando o Laboratório
Computacional}. O objetivo é claro: garantir que, ao virar a última página do
Capítulo~1, você tenha um ambiente \python{} funcional, domine os elementos da
linguagem indispensáveis à química quântica e esteja pronto para instalar e
executar o \pyscf{}. Nada será pressuposto; tudo será construído.

\begin{flushright}
\itshape Cristiano Alves
\end{flushright}
